% ============================================================================
%  บทคัดย่อ / Abstract
% ============================================================================

\chapter*{บทคัดย่อ}
\addcontentsline{toc}{chapter}{บทคัดย่อ}

รายงานฉบับนี้นำเสนอการออกแบบและพัฒนาหุ่นยนต์หยิบและวางชิ้นงานแบบวงกลม (Circular Pick and Place Robot)
สำหรับใช้งานในกระบวนการผลิตเชิงอุตสาหกรรม โดยมีหน้าที่ลำเลียงชิ้นงาน Connection Rod
ระหว่างสถานีทำงานสี่สถานีในลักษณะการหมุนแบบหนึ่งองศาอิสระ

ระบบประกอบด้วยแขนกลที่ขับเคลื่อนด้วยมอเตอร์กระแสตรง RS-775 กำลัง 24~V
ผ่านชุดเฟืองดาวเคราะห์และสายพานฟันที่มีอัตราทดรวม $N = 70$ และประสิทธิภาพ $\eta = 0.836$
ควบคุมด้วยไมโครคอนโทรลเลอร์ STM32G474RE
การออกแบบระบบควบคุมใช้โครงสร้าง Cascade PID สองชั้น ได้แก่
วงรอบความเร็วที่ความถี่ 1~kHz และวงรอบตำแหน่งที่ความถี่ 100~Hz
โดยอ้างอิงพารามิเตอร์ของระบบจากการระบุระบบแบบ Black-Box
ด้วยวิธี Least-Squares ร่วมกับการกรองสัญญาณ Butterworth อันดับสี่ที่ความถี่ตัด 10~Hz
ผลการระบุระบบให้ค่า $R_a = \SI{1.4534}{\ohm}$, $K_t = \SI{0.04065}{\newton\metre\per\ampere}$,
$J = \SI{0.72762}{\kilogram\metre\squared}$ และ $B = \SI{0.19279}{\newton\metre\second\per\radian}$

เพื่อให้บรรลุข้อกำหนดด้าน Cycle Time ไม่เกิน 35~วินาทีสำหรับสี่ชิ้น
ระบบติดตั้ง ZV Input Shaper ที่ออกแบบจากความถี่ธรรมชาติของแท่ง Connection Rod
($\omega_{n,\mathrm{rod}} = \SI{12.386}{\radian\per\second}$, $\zeta_\mathrm{rod} = 0.041$)
ซึ่งช่วยยกเลิกการแกว่งของปลายแขนโดยแทบไม่เพิ่มเวลาหน่วง
นอกจากนี้ยังใช้ Kalman Filter สี่สถานะ
$\mathbf{x} = [\theta_\mathrm{arm},\; \omega_\mathrm{arm},\; i_a,\; \tau_d]^\top$
สำหรับประมาณค่าแรงบิดรบกวนแบบ Real-time
การสื่อสารกับระบบควบคุมส่วนกลางใช้โปรโตคอล Modbus RTU ผ่านบัส RS-485
ส่วนกริปเปอร์ควบคุมด้วยสัญญาณ Digital I/O โดยตรง

ผลการจำลองและการทดสอบระบบย่อยแสดงให้เห็นว่าระบบมีศักยภาพในการบรรลุข้อกำหนดสมรรถนะหลัก
ได้แก่ Settling Time ไม่เกิน 0.5~s (P5), Overshoot ไม่เกิน 1\% (P6),
ความแม่นยำตำแหน่ง $\pm$0.100$\degree$ (A1),
ความซ้ำซากได้ $\pm$0.143$\degree$ ใน 100 รอบ (A2)
และ Cycle time รวม 32.2~วินาที (P1)
ทั้งนี้การทวนสอบข้อกำหนดเชิงปริมาณครบถ้วนอยู่ระหว่างดำเนินการเก็บข้อมูลบนฮาร์ดแวร์จริง

\vspace{1.2cm}
\noindent\textbf{คำสำคัญ:}
หุ่นยนต์หยิบและวางชิ้นงาน,
การควบคุมแบบ Cascade PID,
ZV Input Shaping,
Kalman Filter,
การระบุระบบ,
Modbus RTU

\newpage

% ============================================================================
\chapter*{Abstract}
\addcontentsline{toc}{chapter}{Abstract}

This report presents the design and development of a circular pick-and-place robot
for industrial manufacturing applications, tasked with transferring Connection Rod
workpieces among four work stations via a single rotational degree of freedom.

The mechanical system employs an RS-775 brushed DC motor at 24~V,
coupled through a planetary gearbox and toothed belt drive with a combined gear ratio
of $N = 70$ and efficiency $\eta = 0.836$, controlled by an STM32G474RE microcontroller.
The control architecture adopts a two-loop Cascade PID structure comprising
a velocity inner loop at 1~kHz and a position outer loop at 100~Hz.
Plant parameters were obtained via Black-Box system identification using Least-Squares estimation
with fourth-order Butterworth pre-filtering at 10~Hz, yielding
$R_a = \SI{1.4534}{\ohm}$, $K_t = \SI{0.04065}{\newton\metre\per\ampere}$,
$J = \SI{0.72762}{\kilogram\metre\squared}$, and $B = \SI{0.19279}{\newton\metre\second\per\radian}$.

To satisfy the cycle time requirement of 35~s for four workpieces,
a ZV Input Shaper was designed based on the measured natural frequency of the Connection Rod
($\omega_{n,\mathrm{rod}} = \SI{12.386}{\radian\per\second}$, $\zeta_\mathrm{rod} = 0.041$),
effectively eliminating residual vibration with negligible additional delay.
A four-state Kalman filter with state vector
$\mathbf{x} = [\theta_\mathrm{arm},\; \omega_\mathrm{arm},\; i_a,\; \tau_d]^\top$
provides real-time disturbance torque estimation.
System integration with a central supervisory controller is achieved via
Modbus RTU over RS-485, while the gripper is actuated through digital I/O signals.

Simulation results and subsystem-level testing indicate that the system has the capability
to meet its primary performance requirements,
including settling time within 0.5~s (P5), overshoot not exceeding 1\% (P6),
position accuracy of $\pm$0.100$\degree$ (A1),
repeatability of $\pm$0.143$\degree$ over 100 cycles (A2),
and a total cycle time of 32.2~s for four workpieces (P1).
Full quantitative verification against all requirements is pending hardware-level testing.

\vspace{1.2cm}
\noindent\textbf{Keywords:}
Pick-and-place robot,
Cascade PID control,
ZV input shaping,
Kalman filter,
System identification,
Modbus RTU